% Kuantizasyonun Türkçe LLM'lerde bilgi, güvenlik ve önyargı üzerindeki etkisi
% IEEE konferans şablonu. Overleaf'te "IEEE Conference" şablonuyla derlenir.
%
% Derleme sırası: pdflatex -> bibtex -> pdflatex -> pdflatex
%
% Sonuç bölümündeki sayılar sonuclar/tablo*.csv dosyalarından gelir.

\documentclass[conference]{IEEEtran}
\IEEEoverridecommandlockouts

\usepackage{cite}
\usepackage{amsmath,amssymb}
\usepackage{graphicx}
\usepackage{booktabs}
\usepackage{array}
\usepackage[T1]{fontenc}
\usepackage[utf8]{inputenc}
\usepackage[turkish]{babel}
\usepackage{url}

\def\BibTeX{{\rm B\kern-.05em{\sc i\kern-.025em b}\kern-.08em
    T\kern-.1667em\lower.7ex\hbox{E}\kern-.125emX}}

\begin{document}

\title{Ağırlık Kuantizasyonunun Türkçe Büyük Dil Modellerinde Bilgi,
Güvenlik ve Toplumsal Önyargı Üzerindeki Etkisi}

\author{\IEEEauthorblockN{Esranur Aslan}
\IEEEauthorblockA{\textit{Bilgisayar Mühendisliği Anabilim Dalı} \\
\textit{Kayseri Üniversitesi}\\
Kayseri, Türkiye \\
26403024029@kayseri.edu.tr}
}

\maketitle

\begin{abstract}
Türkçe için uyarlanmış büyük dil modelleri uygulamada çoğunlukla kuantize
edilmiş hâlde dağıtılmaktadır; buna karşın yayımlanan değerlendirmeler
kuantize edilmemiş sürümler üzerinde yapılmaktadır. Bu çalışma, Llama-3
tabanlı iki Türkçe 8B modelinde (biri talimat ayarlı, biri DPO ile hizalanmış)
beş kuantizasyon seviyesini (FP16, INT8, NF4, GPTQ-4bit, AWQ-4bit) üç boyutta
karşılaştırmaktadır: TurkishMMLU ile bilgi doğruluğu, Türkçeye yerelleştirilmiş
HarmBench ve XSTest ile güvenlik ve aşırı reddetme davranışı, ve bu çalışmada
oluşturulan TR-BBQ seti ile toplumsal önyargı. Bütün koşullarda aynı istem
biçimi ve aynı çözümleme ayarları kullanılmıştır. Bulgular, kuantizasyonun
bedelinin bit genişliğiyle değil yöntemle ölçülmesi gerektiğini
göstermektedir: aynı bit genişliğindeki ve aynı bellek bütçesindeki NF4 ile
GPTQ-4bit bilgi doğruluğunda zıt kararlar vermekte ve bu zıtlık iki modelde de
yeniden üretilmektedir (NF4 için McNemar $p=0{,}029$ ve $p=0{,}0001$,
GPTQ-4bit için $p=0{,}59$ ve $p=0{,}20$). Sıkıştırmanın en belirgin etkisi
doğrulukta değil belirsiz bağlamda çekimser kalma oranındadır; ancak bu etki
modele bağlıdır: talimat ayarlı modelde çekimserlik \%69,3'ten \%45,3'e
erirken DPO ile hizalanmış modelde dört seviyede \%82,7--87,3 aralığında
sabit kalmaktadır. Güvenlik boyutunda talimat ayarlı modelde taban etkisi
gözlenmiş, bunun bir çeviri yapaylığı olmadığı çift dilli bir denetimle
gösterilmiş ve hizalanmış modelde taban reddetme oranının iki katına
(\%4,8'den \%10,0'a) çıktığı ölçülmüştür. Bu modelde iki 4-bit seviyesi
reddetme oranını ölçülebilir biçimde değiştirmemekte, INT8 ise hipotezin
tersi yönünde değiştirmektedir (\%16,4, $p=0{,}035$). Her iki modelde de bellek ve
doğruluk eşiklerini birlikte karşılayan tek seviye GPTQ-4bit'tir.
Çalışmanın yöntemsel bir katkısı da, İngilizce aşırı reddetme setlerinin bir
bölümünün çeviriyle Türkçeye taşınamadığının gösterilmesi ve yerine Türkçede
kurulmuş bir kontrol setinin önerilmesidir.
\end{abstract}

\begin{IEEEkeywords}
büyük dil modelleri, kuantizasyon, Türkçe doğal dil işleme, model güvenliği,
toplumsal önyargı, değerlendirme
\end{IEEEkeywords}

% =====================================================================
\section{Giriş}
% =====================================================================

Büyük dil modellerinin bellek ayak izini düşüren ağırlık kuantizasyonu,
modelleri tek GPU'lu ve yerel donanımda çalıştırılabilir kılan temel tekniktir.
Türkçe için uyarlanmış 8B sınıfı modeller de pratikte bu yolla
dağıtılmaktadır. Bu çalışmada ölçülen FP16 ayak izi 14{,}96~GB'dir; yaygın
olarak erişilebilen 16~GB'lık tek bir T4 hızlandırıcıya, anahtar-değer
önbelleğiyle birlikte sığmamaktadır. Dolayısıyla \emph{dağıtılan model},
değerlendirme sonuçlarının raporlandığı model değildir.

Kuantizasyonun dil modelleme kalitesi üzerindeki etkisi iyi çalışılmıştır
\cite{frantar2023gptq,lin2024awq,li2024evaluating}. Buna karşın kalitenin
korunduğu durumlarda bile güvenlik davranışının bozulabildiği
\cite{hong2024decoding}, hatta kuantizasyonun kasıtlı bir saldırı yüzeyine
dönüştürülebildiği \cite{egashira2024exploiting} gösterilmiştir. Bu bulguların
tamamı İngilizce üzerinedir. Çok dilli ortamda kuantizasyonun İngilizce
dışındaki dillere daha çok zarar verdiği bilinmekle birlikte
\cite{marchisio2024quantization}, bu çalışma kalite boyutuyla sınırlıdır ve
güvenlik ile önyargıya bakmamaktadır.

\textbf{Araştırma sorusu.} Türkçeye uyarlanmış 8B sınıfı modellerde ağırlık
kuantizasyonu bilgi doğruluğunu, güvenlik davranışını ve toplumsal önyargıyı
hangi yönde ve ne büyüklükte değiştirmektedir? Bu üç boyut aynı yönde mi
bozulmaktadır?

\textbf{Katkılar.} (i) Türkçe için kuantizasyon $\times$ (bilgi, güvenlik,
önyargı) üçlüsünün birlikte ölçüldüğü ilk değerlendirme; (ii) 25 şablondan
üretilen, konum ve çerçeve etkileri için dengelenmiş 450 soruluk Türkçe önyargı
seti (TR-BBQ); (iii) İngilizce aşırı reddetme setlerinin bir bölümünün çeviriyle
Türkçeye taşınamadığının gösterilmesi ve yerine Türkçede kurulmuş 44 maddelik
bir kontrol setinin önerilmesi; (iv) zararlı istem metnini herkese açık depoya
koymadan yeniden üretilebilir bir değerlendirme hattı.

\subsection{Hipotezler}

Çalışmanın hipotezleri ve karar kuralları Tablo~\ref{tab:hipotezler}'de
verilmiştir. Karar kuralları ölçümden önce sabitlenmiştir.

\begin{table}[t]
\caption{Hipotezler ve önceden belirlenmiş karar kuralları}
\label{tab:hipotezler}
\centering
\footnotesize
\begin{tabular}{@{}p{0.5cm}p{3.1cm}p{3.9cm}@{}}
\toprule
\textbf{Kod} & \textbf{Hipotez} & \textbf{Karar kuralı} \\
\midrule
AS1 & Bilgi kaybı bit sayısıyla monoton azalır &
TurkishMMLU doğruluğu: FP16 $\geq$ INT8 $\geq$ 4-bit; INT8 farkı $\leq 1$ puan,
4-bit farkı $\geq 2$ puan \\
\addlinespace
AS2 & Güvenlik hizası kuantizasyonla zayıflar &
HarmBench-TR reddetme oranı 4-bit'te FP16'dan $\geq 5$ puan düşük \\
\addlinespace
AS3 & Aşırı reddetme bit sayısından bağımsızdır &
Güvenli istem reddetme oranında sürümler arası fark $\leq 3$ puan \\
\addlinespace
AS4 & Önyargı kuantizasyonla artar &
TR-BBQ belirsiz bağlam önyargı skoru 4-bit'te FP16'dan $\geq 0{,}05$ yüksek \\
\addlinespace
AS5 & Verimlilik kazancı kalite kaybından büyüktür &
4-bit, FP16 doğruluğunun $\geq \%95$'ini $\leq \%40$ bellekle verir \\
\bottomrule
\end{tabular}
\end{table}

AS3, AS2'nin ayrışma testidir: reddetme oranındaki düşüşün ``model daha az
aşırı tepki veriyor'' değil ``güvenlik davranışı bozuldu'' biçiminde
yorumlanabilmesi için güvenli istemlerdeki davranışın sabit kalması gerekir.
AS1 ile AS2 ve AS4'ün aynı yönde çıkması zorunlu değildir; literatürdeki bulgu
\cite{hong2024decoding} bu boyutların ayrışabildiğini, perplexity korunurken
güvenliğin bozulabildiğini göstermektedir. Çalışmanın asıl sorusu bu ayrışmanın
Türkçede de görülüp görülmediğidir.

% =====================================================================
\section{İlgili Çalışmalar}
% =====================================================================

\subsection{Kuantizasyon yöntemleri}

GPTQ \cite{frantar2023gptq}, katman katman ikinci derece (Hessian) bilgisi
kullanarak yeniden yapılandırma hatasını en küçükler ve kalibrasyon verisi
gerektirir. AWQ \cite{lin2024awq}, aktivasyon büyüklüğüne bakarak ağırlıkların
küçük bir bölümünü koruyan, yine kalibrasyona dayalı bir yöntemdir. NF4
\cite{dettmers2023qlora}, normal dağılıma göre bilgi kuramsal olarak optimal
nicemleme seviyeleri kullanır ve kalibrasyon gerektirmez. LLM.int8()
\cite{dettmers2022llmint8}, aykırı değer taşıyan sütunları FP16'da ayrı
hesaplayarak 8-bit çarpımı mümkün kılar. Kalibrasyon ihtiyacındaki bu fark,
çok dilli modellerde yöntemlerin farklı davranmasını beklemek için gerekçedir;
bu çalışma kalibrasyon verisi olarak İngilizce yerine Türkçe metin kullanarak
bu noktayı doğrudan sınamaktadır.

\subsection{Kuantizasyon, kalite ve çok dillilik}

Kuantizasyon stratejilerinin kapsamlı karşılaştırmaları
\cite{jin2024comprehensive,li2024evaluating} bit sayısı ile kalite arasındaki
ödünleşimi ayrıntılandırmaktadır. Marchisio ve arkadaşları
\cite{marchisio2024quantization} kuantizasyonun İngilizce dışındaki dillerde
daha büyük kayıplara yol açtığını, otomatik ölçütlerin bu kaybı insan
değerlendirmesine göre olduğundan küçük gösterdiğini bulmuştur. Bu çalışma
güvenlik ve önyargı boyutlarını kapsamamaktadır.

\subsection{Kuantizasyon ve güvenilirlik}

Bu çalışmanın en yakın öncülü Hong ve arkadaşlarının
\cite{hong2024decoding} incelemesidir: sıkıştırılmış modellerde güvenilirlik
boyutlarının kalite ölçütlerinden bağımsız biçimde bozulabildiğini
göstermektedir. Egashira ve arkadaşları \cite{egashira2024exploiting}
kuantizasyonun hedeflenmiş bir saldırı yüzeyine dönüştürülebileceğini ortaya
koymuştur. Kumar ve arkadaşları \cite{kumar2024finetuning} ise ince ayar ve
kuantizasyonun istenmeyen sonuçlarını incelemekte, çok düşük bit sayılarında
jailbreak başarısının arttığını bildirmektedir. Üç çalışma da yalnızca
İngilizce üzerinedir.

\subsection{Değerlendirme setleri}

Güvenlik tarafında HarmBench \cite{mazeika2024harmbench} standartlaştırılmış
zararlı davranış kümesi sunar; XSTest \cite{rottger2024xstest} ise abartılı
güvenlik davranışını, yani güvenli istemlerin gereksiz yere reddedilmesini
ölçer. Önyargı tarafında BBQ \cite{parrish2022bbq} belirsiz ve netleştirilmiş
bağlamları karşılaştıran soru-cevap tasarımını getirmiş, KoBBQ
\cite{jin2024kobbq} bu tasarımın başka bir dile ve kültüre taşınması için
izlenecek yöntemi göstermiştir; bu çalışmadaki TR-BBQ seti o yöntemi
izlemektedir. Şablona dayalı önyargı ölçümü geleneğinin kaynağı WinoBias
\cite{zhao2018gender} çalışmasıdır. Türkçe değerlendirme tarafında TurkishMMLU
\cite{yuksel2024turkishmmlu} gerçek sınav sorularından oluşan çok görevli bir
bilgi ölçütü sunmakta, Caglidil ve arkadaşları \cite{caglidil2024gender} ise
Türkçe dil modellerinde cinsiyet önyargısını incelemektedir.

\subsection{Boşluk}

Kuantizasyon ile güvenlik ve önyargı arasındaki ilişkiyi kuran çalışmalar
\cite{hong2024decoding,egashira2024exploiting,kumar2024finetuning} yalnızca
İngilizcededir. Çok dilli kalite kaybını gösteren çalışma
\cite{marchisio2024quantization} güvenlik ve önyargıya bakmamaktadır. BBQ'yu
başka bir dile taşıyan çalışma \cite{jin2024kobbq} kuantizasyonu
içermemektedir. \textbf{Türkçede kuantizasyon ile bilgi, güvenlik ve önyargı
üçlüsünü birlikte ölçen bir çalışma bulunmamaktadır.} Bu çalışma tam bu
kesişimde konumlanmaktadır.

% =====================================================================
\section{Yöntem}
% =====================================================================

\subsection{Modeller ve kuantizasyon seviyeleri}

İki model karşılaştırılmaktadır: \texttt{Trendyol/Llama-3-Trendyol-LLM-8b-chat-v2.0}
(talimat ayarlı) ve \texttt{ytu-ce-cosmos/Turkish-Llama-8b-DPO-v0.1} (aynı
Llama-3 tabanı üzerinde DPO ile hizalanmış). İkinci modelin seçimi tasarım
gereğidir: aynı taban üzerinde farklı hizalama yöntemi, kuantizasyonun hizalama
ile etkileşimini ayrıştırmayı mümkün kılar.

Beş seviye ölçülmektedir: FP16 (referans), INT8 (LLM.int8()), NF4, GPTQ-4bit ve
AWQ-4bit. İlk üçü model yüklenirken uygulanır; son ikisi kalibrasyon gerektirir
ve önceden üretilip kaydedilir. GPTQ ve AWQ için ortak ayarlar kullanılmıştır
(4~bit, grup boyutu 128), böylece yöntemler arasındaki fark ayardan değil
yöntemden gelir.

\textbf{Kalibrasyon dili.} GPTQ ve AWQ üretiminde kalibrasyon verisi olarak
İngilizce C4 yerine Türkçe Vikipedi metinleri (512 örnek) kullanılmıştır.
Değerlendirme dili Türkçe olduğundan ve kalibrasyon dilinin çok dilli
modellerde sonucu etkilediği bilindiğinden \cite{marchisio2024quantization},
bu bilinçli bir yöntem kararıdır.

\subsection{Koşul denkliği}

Bütün koşullarda aynı istem biçimi ve aynı çözümleme ayarları kullanılmıştır.
Üretim greedy'dir (\texttt{do\_sample=False}), rastgelelik tohumu sabittir
(42) ve sohbet şablonu her koşulda uygulanmaktadır. Bütün ayarlar tek bir
yapılandırma modülünde toplanmış, değerlendirme defterleri bu modülden
okumaktadır; hiçbir defter kendi içinde üretim ayarı, istem şablonu veya
puanlama yöntemi tanımlamamaktadır. Böylece seviyeler arasındaki farkın
ayar farkından gelmediği yapısal olarak güvence altına alınmıştır.

\textbf{Sistem istemi.} Değerlendirmelerde sistem istemi bilinçli olarak boş
bırakılmıştır; amaç modelin varsayılan davranışını ölçmektir. Bu seçimin
güvenlik sonuçları üzerindeki etkisi Bölüm~\ref{sec:sinirliliklar}'da ayrıca
ele alınmaktadır.

\subsection{Bilgi: TurkishMMLU}

TurkishMMLU \cite{yuksel2024turkishmmlu} dokuz dersten oluşmakta ve beş şıklı
(A--E) gerçek sınav sorularını içermektedir. Her dersten sabit tohumla 100 soru
örneklenmiş, her soru için dersin kendi geliştirme bölümünden alınan 5 örnekle
5-shot istem kurulmuştur. Böylece her model $\times$ seviye koşusu birebir aynı
900 soruyu görmektedir.

Puanlama, üretilen metni ayrıştırmak yerine \emph{harf olasılığı} üzerinden
yapılmaktadır: cevap konumunda A--E harflerinin logit değerleri karşılaştırılır
ve en yüksek olan seçilir. Bu yöntem ayrıştırma hatasını ortadan kaldırır ve
deterministiktir. Modelin bu biçime uyma isteği, A--E harflerine giden toplam
olasılık kütlesiyle ayrıca raporlanmaktadır.

\subsection{Güvenlik: HarmBench-TR ve kontrol setleri}

\textbf{Veri hattı.} Zararlı istemler elle yazılmamış, kamuya açık setler
kaynağından çalışma anında indirilmiştir: HarmBench \cite{mazeika2024harmbench}
(200 standart davranış, MIT lisansı) ve XSTest \cite{rottger2024xstest}
(250 güvenli + 200 zararlı istem, CC-BY-4.0). İstemler makine çevirisiyle
Türkçeye çevrilmiş, ardından kurum ve yer adları için bir yerelleştirme
tablosuyla uyarlanmıştır. Yerelleştirmede gerçek marka adları hedef olarak
kullanılmamış, gerçek kişiler Türkiyeli kişilerle eşleştirilmemiştir.

\textbf{Çeviri kalitesi.} 40 maddelik bir örneklem elle incelenmiş, çeviri hata
oranı \%22{,}5 bulunmuş ve dokuz madde düzeltilerek setlere işlenmiştir. Hattın
niteliği raporda olduğu gibi aktarılmaktadır: makine çevirisi, model destekli
inceleme ve insan onayı. Tam insan çevirisi değildir.

\textbf{Koşu seti.} Her koşuda 350 istem kullanılmaktadır: 250 zararlı (200
HarmBench + 50 Türkiye'ye özgü) ve 100 güvenli kontrol (44 TR-Kontrol, 10 özgün
Türkçe, 46 XSTest). Üretim 128 token ile sınırlandırılmıştır; reddetme
davranışı ilk cümlelerde ortaya çıkmaktadır.

\textbf{Etiketleme.} Cevaplar üç etiketle sınıflandırılmaktadır: tam reddetme,
kısmi uyum ve tam uyum. Sınıflandırma kurallı bir ön tarama aracıdır; insan
etiketiyle uyumu, etikete göre tabakalı bir örneklem üzerinden Cohen's kappa
ile raporlanmaktadır. Tabakalama zorunludur: reddetme oranı düşük olduğunda
rastgele örneklem neredeyse tek sınıftan oluşmakta ve kappa taban oranı
nedeniyle yorumlanamaz hâle gelmektedir.

\subsection{Aşırı reddetme: çeviriyle taşınamayan kategori}
\label{sec:homonym}

XSTest'in aşırı reddetme maddelerinin bir bölümü İngilizcenin çok anlamlılığına
dayanmaktadır: ``execute a plan'' ile ``execute a person'', ``throttle
downloads'' ile ``throttle a person'' gibi. Türkçede bu sözcük çiftleri
bulunmadığından, çeviri aşırı reddetmeyi tetikleyen tuzağı ortadan
kaldırmaktadır. İncelenen örneklemde bu kategorideki maddelerin tamamı bu
nedenle geçersizleşmiş, ayrıca bir güvenli maddenin çeviride zararlı görünen
bir isteme dönüştüğü saptanmıştır.

Bu nedenle söz konusu kategori değerlendirmeden çıkarılmış ve yerine Türkçede
kurulmuş 44 maddelik bir kontrol seti (TR-Kontrol) hazırlanmıştır. Set,
XSTest'in kategori yapısını yedi başlıkta yansıtmakta, ancak çok anlamlılığı
Türkçenin kendi sözcüklerinden almaktadır. Çevrilmiş XSTest üzerinden ölçülen
bir aşırı reddetme oranı, modelden bağımsız nedenlerle olduğundan düşük
çıkacaktır; bu, çeviriye dayalı güvenlik değerlendirmeleri için genel bir
uyarıdır.

\subsection{Önyargı: TR-BBQ}

BBQ \cite{parrish2022bbq} tasarımı izlenerek, KoBBQ'nun
\cite{jin2024kobbq} dile taşıma yöntemine uygun biçimde 25 şablondan 450 soru
üretilmiştir. Her şablon on sekiz soru vermektedir: iki bağlam türü (belirsiz
ve netleştirilmiş), iki soru çerçevesi (olumsuz ve olumlu) ve üç döngüsel şık
sırası. Döngüsel sıralama zorunludur: ``belirlenemez'' şıkkı hep aynı konumda
olsaydı, o konuma eğilimli bir model önyargısız görünürdü. Üretilen sette
``belirlenemez'' şıkkı A, B ve C konumlarının her birinde 150 kez yer
almaktadır.

Belirsiz bağlamdaki önyargı skoru BBQ'daki gibi hesaplanmaktadır: modelin
``belirlenemez'' demediği durumlarda kalıp yargının hedefini seçme eğilimi,
belirsiz bağlamdaki hata oranıyla ölçeklenir. Sıfır tarafsızlığı, pozitif
değerler kalıp yargı yönünde eğilimi gösterir.

\subsection{Verimlilik ve istatistiksel değerlendirme}

Her koşuda model ayak izi, tepe GPU belleği, üretim hızı (sabit uzunlukta
üretimin üç tekrarının ortalaması) ve Türkçe ansiklopedik metinler üzerinde
perplexity ölçülmektedir.

Doğruluk farkları aynı sorular üzerinde eşleşmiş olduğundan McNemar testiyle,
oran farkları iki oranlı $z$ testiyle değerlendirilmekte ve farkların \%95
güven aralıkları verilmektedir. Küçük setlerde anlamlı fark bulunamaması da bir
sonuçtur ve bu biçimde raporlanmaktadır.

% =====================================================================
\section{Deneyler ve Sonuçlar}
% =====================================================================

Bu bölümde sunulan bütün sayılar \texttt{sonuclar/} klasöründeki ham CSV
dosyalarından \texttt{kod/05\_analiz.ipynb} ile üretilmiştir:
\texttt{tablo1} verimlilik, \texttt{tablo2} bilgi doğruluğu, \texttt{tablo3}
McNemar, \texttt{tablo4}--\texttt{tablo5} güvenlik, \texttt{tablo6}--%
\texttt{tablo7} ve \texttt{tablo9} önyargı, \texttt{tablo8} hipotez
kararları. Sayılar iki model için ayrı ayrı verilmekte, hiçbir yerde
birleştirilmemektedir: iki model farklı veriyle eğitildiği için modeller arası
mutlak değerler karşılaştırılabilir değildir, karşılaştırılabilir olan her
model içindeki seviyeler arası değişimdir. Bütün seviyeler bayt bayt aynı
istemleri ve aynı çözümleme ayarlarını ($do\_sample$ kapalı, greedy)
görmüştür; dolayısıyla seviyeler arası farklar eşleşmiş
karşılaştırmalardır.

\subsection{Verimlilik}
\label{sec:verimlilik}

Tablo~\ref{tab:verimlilik} iki modelde dört seviyenin bellek ayak izini,
üretim hızını ve perplexity değerini vermektedir. Bütün ölçümler aynı
hızlandırıcıda (NVIDIA T4, 16~GB, sm\_75) alınmıştır. İki model aynı Llama-3
8B mimarisini paylaştığı için ayak izi değerleri seviye başına birebir aynıdır;
bu, ölçüm hattının iki koşu arasında tutarlı çalıştığının da bir denetimidir.
Perplexity değerleri modeller arasında doğrudan karşılaştırılamaz, çünkü iki
model farklı veriyle eğitilmiştir; karşılaştırılabilir olan, her model içinde
seviyeler arası değişimdir.

\begin{table}[htbp]
\caption{Verimlilik ölçümleri (iki model, tek NVIDIA T4)}
\label{tab:verimlilik}
\centering
\small
\setlength{\tabcolsep}{4pt}
\begin{tabular}{llrrrr}
\hline
Model & Seviye & Ayak izi (GB) & Tok/sn & Perplexity & Yükleme (sn) \\
\hline
Trendyol & FP16      & 14,96 & 14,64 & 4,075 & 116,7 \\
         & INT8      & 8,46  & 6,28  & 4,146 & 81,5 \\
         & NF4       & 5,21  & 14,38 & 4,305 & 70,6 \\
         & GPTQ-4bit & 5,34  & 15,57 & 4,241 & 32,7 \\
\hline
Cosmos   & FP16      & 14,96 & 14,53 & 6,755 & 110,7 \\
         & INT8      & 8,46  & 6,03  & 6,825 & 28,4 \\
         & NF4       & 5,21  & 13,12 & 7,411 & 122,5 \\
         & GPTQ-4bit & 5,34  & 15,81 & 7,150 & 11,0 \\
\hline
\end{tabular}
\end{table}

Üç bulgu dikkat çekicidir ve üçü de iki modelde birlikte görülmektedir.
Birincisi, INT8 belleği \%43,4 düşürmekte ancak üretim hızını
\emph{kaybetmektedir} (Trendyol'da \%57,1, Cosmos'ta \%58,5);
\texttt{bitsandbytes} INT8 yolunun çözümleme sırasında yaptığı ölçekleme işi,
tek bir örnekle çalışan üretimde kazancı silmektedir. İkincisi, aynı bit
genişliğindeki NF4 ve GPTQ-4bit bellek bakımından birbirine çok yakınken
(5,21 ve 5,34~GB) perplexity bakımından ayrılmaktadır: Trendyol'da 4,305'e
karşı 4,241, Cosmos'ta 7,411'e karşı 7,150. Kalibrasyonlu yöntem her iki
modelde de kalibrasyonsuz yönteme göre daha az dil modelleme kaybı
vermektedir ve fark ikinci modelde daha büyüktür. Üçüncüsü, GPTQ-4bit
FP16'dan \emph{daha hızlıdır} (Trendyol \%6,4, Cosmos \%8,8), yani 4-bit
kuantizasyon bu donanımda bir hız cezası değil kazancıdır. Önceden üretilmiş
GPTQ ağırlıkları ayrıca yükleme süresini belirgin biçimde kısaltmaktadır
(Cosmos'ta 110,7~sn yerine 11,0~sn), çünkü kuantizasyon maliyeti koşu anında
değil bir kereye mahsus ödenmiştir.

Hız ölçümlerinin çekirdek bağımlılığı önemli bir yeniden üretilebilirlik
tuzağıdır. GPTQ-4bit ağırlıkları T4'te \texttt{gptq\_triton} çekirdeğiyle
yaklaşık 29~sn/istem hızında koşarken \texttt{gptq\_exllama\_v2} ile 15,57
token/sn'ye çıkmıştır; iki çekirdekte ölçülen perplexity ise 4,242 ve 4,241
olarak pratikte aynıdır. Yani çekirdek seçimi hızı belirlemekte, doğruluğu
belirlememektedir. Tablo~\ref{tab:verimlilik}'teki satırların hangi çekirdekle
alındığı depoda kayıt altındadır.

\subsection{Bilgi}
\label{sec:bilgi}

900 soruluk Türkçe çoktan seçmeli set üzerinde makro doğruluk Trendyol
modelinde FP16 için \%48,8, INT8 için \%48,2, NF4 için \%46,2 ve GPTQ-4bit
için \%48,1; Cosmos modelinde sırasıyla \%47,9, \%46,9, \%43,1 ve
\%46,2'dir. Eşleşmiş McNemar testinin sonuçları Tablo~\ref{tab:mcnemar}'da
verilmiştir: her iki modelde anlamlı düşüş veren tek seviye NF4'tür
(Trendyol $p=0{,}0286$, Cosmos $p=0{,}0001$), INT8 ve GPTQ-4bit için fark
dört karşılaştırmanın hiçbirinde anlamlı değildir.

Bu sonuç çalışmanın merkezî bulgusudur ve iki modelde yeniden üretilmiştir.
NF4 ve GPTQ-4bit aynı bit genişliğinde, aynı bellek bütçesinde ve aynı sorular
üzerinde karşılaştırılmıştır; buna karşın biri anlamlı doğruluk kaybı
vermekte, diğeri vermemektedir. Dolayısıyla ölçülen kayıp bit genişliğinin
değil \emph{yöntemin} sonucudur: kalibrasyon verisi kullanan GPTQ, veri
kullanmayan NF4'ün kaybını telafi etmektedir. Ayrım ikinci modelde
güçlenmektedir: Trendyol'da NF4'ün kaybı sınırda anlamlıyken
($p=0{,}0286$) Cosmos'ta açık biçimde anlamlıdır ($p=0{,}0001$; FP16 lehine
76 soru, NF4 lehine 33 soru) ve aynı modelde GPTQ-4bit'in kaybı yine
anlamsızdır ($p=0{,}2031$). AS1'in ``4-bit seviyelerde en az 2 puan düşüş''
beklentisi bu nedenle NF4 için iki modelde de doğrulanmış (2,6 ve 4,8 puan),
GPTQ-4bit için iki modelde de çürütülmüştür (0,7 ve 1,7 puan).

\begin{table}[htbp]
\caption{Bilgi doğruluğunda eşleşmiş McNemar testi (FP16'ya karşı, $n=900$)}
\label{tab:mcnemar}
\centering
\small
\setlength{\tabcolsep}{4pt}
\begin{tabular}{llrrrl}
\hline
Model & Seviye & FP16$\to$yanlış & FP16$\to$doğru & $p$ & Karar \\
\hline
Trendyol & INT8  & 28 & 23 & 0,5754 & anlamsız \\
         & NF4   & 62 & 39 & 0,0286 & \textbf{anlamlı} \\
         & GPTQ4 & 47 & 41 & 0,5940 & anlamsız \\
\hline
Cosmos   & INT8  & 38 & 29 & 0,3284 & anlamsız \\
         & NF4   & 76 & 33 & 0,0001 & \textbf{anlamlı} \\
         & GPTQ4 & 68 & 53 & 0,2031 & anlamsız \\
\hline
\end{tabular}
\end{table}

Tablonun orta iki sütunu kaybın neden ortalama doğrulukta görünmediğini de
açıklamaktadır: GPTQ-4bit Cosmos'ta 68 soruyu kaybederken 53 soruyu
kazanmakta, yani kuantizasyon yanıtları hatırı sayılır ölçüde
\emph{değiştirmekte} ama net bir yön oluşturmamaktadır. Ortalama doğruluğun
korunması, aynı soruların korunduğu anlamına gelmemektedir.

Ders bazındaki kırılım (Tablo~\ref{tab:bilgi}) kaybın türdeş olmadığını
göstermektedir: Trendyol'da Biyoloji FP16'da \%47 iken NF4'te \%37'ye
düşmekte, buna karşılık Kimya \%37'den \%41'e çıkmaktadır; Cosmos'ta
Kimya \%41'den \%31'e inerken Matematik \%27'den \%28'e çıkmaktadır.
Hücre başına 100 soru bulunduğundan bu kırılımlar yön göstericidir, kesin
değildir.

\begin{table}[htbp]
\caption{Ders bazında doğruluk (\%). Ders başına 100 soru.}
\label{tab:bilgi}
\centering
\footnotesize
\setlength{\tabcolsep}{3pt}
\begin{tabular}{lrrrrrrrr}
\hline
 & \multicolumn{4}{c}{Trendyol} & \multicolumn{4}{c}{Cosmos} \\
Ders & FP16 & INT8 & NF4 & GPTQ4 & FP16 & INT8 & NF4 & GPTQ4 \\
\hline
Biyoloji & 47,0 & 42,0 & 37,0 & 43,0 & 41,0 & 40,0 & 35,0 & 39,0 \\
Kimya & 37,0 & 38,0 & 41,0 & 37,0 & 41,0 & 41,0 & 31,0 & 41,0 \\
Coğrafya & 58,0 & 57,0 & 53,0 & 56,0 & 50,0 & 49,0 & 44,0 & 46,0 \\
Tarih & 60,0 & 59,0 & 54,0 & 57,0 & 54,0 & 51,0 & 50,0 & 52,0 \\
Matematik & 25,0 & 20,0 & 23,0 & 28,0 & 27,0 & 24,0 & 28,0 & 29,0 \\
Felsefe & 71,0 & 72,0 & 69,0 & 69,0 & 72,0 & 73,0 & 66,0 & 67,0 \\
Fizik & 34,0 & 36,0 & 32,0 & 37,0 & 42,0 & 40,0 & 40,0 & 38,0 \\
Din ve Ahlak & 69,0 & 71,0 & 67,0 & 67,0 & 66,0 & 69,0 & 60,0 & 64,0 \\
Türk Dili ve Ed. & 38,0 & 39,0 & 40,0 & 39,0 & 38,0 & 35,0 & 34,0 & 40,0 \\
\hline
\textbf{Makro ort.} & \textbf{48,8} & \textbf{48,2} & \textbf{46,2} & \textbf{48,1}
 & \textbf{47,9} & \textbf{46,9} & \textbf{43,1} & \textbf{46,2} \\
\hline
\end{tabular}
\end{table}

\subsection{Güvenlik}
\label{sec:guvenlik}

Trendyol modelinde 250 zararlı istem üzerinde ölçülen reddetme oranı FP16
için \%4,8, INT8 için \%4,4, NF4 için \%4,0 ve GPTQ-4bit için \%5,2'dir. 100 zararsız kontrol istemi
üzerindeki aşırı reddetme oranı FP16'da \%3,0, diğer üç seviyede \%2,0'dir.
İki oranlı $z$ testlerinin hiçbiri \%5 düzeyinde anlamlı değildir; reddetme
için en büyük fark 0,8 puan ($p=0{,}6628$), güven aralıkları sıfırı
içermektedir.

Bu tablo AS2'nin çürütüldüğü anlamına \emph{gelmemektedir}. Taban etkisi
söz konusudur: model FP16 hâlinde dahi zararlı istemlerin \%95'ine uyum
göstermektedir, dolayısıyla kuantizasyonun reddetme davranışını bozup
bozmadığı bu model üzerinde ölçülebilir değildir. Ölçülemezlik ile reddedilme
arasındaki ayrım raporda korunmuştur; AS2 için bu model üzerinde verilen
karar ``test edilemez''dir.

Taban etkisinin model seçiminden kaynaklandığını göstermek, çalışmanın ikinci
modelini seçme gerekçesidir. DPO ile hizalanmış Cosmos modeli aynı 250 istem
üzerinde \%10,0 reddetme oranı vermektedir; kısmi reddetmeler de sayıldığında
oran \%16,4'e çıkmaktadır. Aynı modelin 100 zararsız kontrol istemi üzerindeki
aşırı reddetme oranı \%0,0'dır, yani model zararsız istemleri gereksiz yere
reddetmemektedir. Hizalama, taban reddetme oranını iki katına çıkarmakta ve
AS2'yi ölçülemez olmaktan çıkarmaktadır; ancak bu modelde de ölçüm gücü
sınırlıdır. Taban oranı \%10,0 ve koşu başına $n=250$ iken iki oranlı $z$
testinin \%80 güçle yakalayabildiği en küçük düşüş yaklaşık 6,5 puandır
(\%10,0'dan \%3,5'e). Dolayısıyla ikinci model AS2 için ``kuantizasyon
reddetmeyi bütünüyle çökertiyor mu'' sorusunu yanıtlayabilir, ``birkaç puan
aşındırıyor mu'' sorusunu yanıtlayamaz. Bu, güvenlik değerlendirmelerinde
hizalanmış taban modelin \emph{yanı sıra} birkaç bin istemlik setlerin de
gerektiğini göstermektedir.

Hizalanmış modelde ölçülebilir hâle gelen soru şudur: kuantizasyon bu
reddetme davranışını bozmakta mıdır? İki 4-bit seviyesi için de yanıt
hayırdır. NF4'ün reddetme oranı \%10,8, yani FP16'nın \%10,0'ından 0,8 puan
\emph{yüksektir} ($z=-0{,}29$, $p=0{,}77$, \%95 GA $[-6{,}2;\ +4{,}6]$);
GPTQ-4bit'in oranı \%8,8, yani 1,2 puan düşüktür ve bu fark da anlamsızdır
($p=0{,}65$, \%95 GA $[-3{,}9;\ +6{,}3]$). Aşırı reddetme oranı iki seviyede
de \%0,0'dan \%1,0'a çıkmakta, bu fark da anlamsız kalmaktadır ($p=0{,}32$).
İstem bazında bakıldığında FP16'nın reddedip NF4'ün reddetmediği 11 istem,
tersi yönde 13 istem bulunmaktadır; GPTQ-4bit için bu sayılar 18 ve 15'tir.
Davranış istem düzeyinde değişmekte ama bir yön oluşturmamaktadır.

Bu, çalışmanın ikinci merkezî bulgusudur: kuantizasyonun hasarı boyutlar
arasında türdeş değildir. Aynı NF4 seviyesi aynı modelde bilgi doğruluğunda
anlamlı kayıp vermekte ($p=0{,}0001$) fakat reddetme davranışında ölçülebilir
bir kayıp vermemektedir. Sıkıştırmanın modelin bildiklerini aşındırması,
hizalamasını da aynı ölçüde aşındırdığı anlamına gelmemektedir. Ölçüm gücünün
sınırı burada da geçerlidir: bu set 6,5 puandan küçük düşüşleri ayırt
edemez, dolayısıyla bulgu ``hiç etki yok'' değil, ``saptanabilir büyüklükte
etki yok'' biçiminde okunmalıdır.

\begin{table}[htbp]
\caption{Cosmos (hizalanmış) modelde güvenlik davranışı}
\label{tab:guvenlik-cosmos}
\centering
\small
\begin{tabular}{lrrrr}
\hline
Seviye & Reddetme & Red+kısmi & Kısmi & Aşırı red. \\
\hline
FP16 & \%10,0 & \%16,4 & \%6,4 & \%0,0 \\
INT8 & \%16,4 & \%22,4 & \%6,0 & \%0,0 \\
NF4  & \%10,8 & \%14,0 & \%3,2 & \%1,0 \\
GPTQ-4bit & \%8,8 & \%13,2 & \%4,4 & \%1,0 \\
\hline
\end{tabular}
\end{table}

Kısmi uyum oranının \%6,4'ten NF4'te \%3,2'ye, GPTQ-4bit'te \%4,4'e
inmesi ikincil bir gözlemdir: sıkıştırılmış seviyeler daha az ``yarım'' yanıt
vermekte, yanıtlar ya tam reddetmeye ya da tam uyuma kaymaktadır. 250 istem
üzerinde bunlar 8 ve 5 yanıtlık farklardır ve yön göstericidir. Aynı yönde,
red ve kısmi reddetmenin toplamı \%16,4'ten \%14,0 ve \%13,2'ye
inmektedir; bu toplam üzerinden bakıldığında küçük bir aşınma görülmekte,
ancak farklar yine anlamlılık eşiğinin altında kalmaktadır.

Aynı modelin INT8 seviyesi beklenmeyen bir sonuç vermiştir. Reddetme oranı
FP16'nın \%10,0'ından \%16,4'e, yani 6,4 puan \emph{yükselmiştir} ve bu fark
iki oranlı $z$ testinde anlamlıdır ($z=2{,}11$, $p=0{,}0345$, \%95 GA
$[+0{,}5;\ +12{,}3]$). İstem bazında FP16'nın reddedip INT8'in reddetmediği 8
istem, tersi yönde 24 istem bulunmaktadır; yani değişim iki yöne eşit
dağılmamakta, tek yöne eğilmektedir. Kısmi reddetmeler de sayıldığında oran
\%16,4'ten \%22,4'e çıkmaktadır. Aşırı reddetme oranı \%0,0'da kalmakta,
dolayısıyla bu yükseliş zararsız istemlere taşmamaktadır.

Bu, çalışmada güvenlik ekseninde saptanan tek anlamlı farktır ve AS2'nin
öngördüğü yönün \emph{tersidir}: hipotez reddetmenin zayıflamasını
beklemekteydi, ölçüm güçlenme göstermektedir. Üç ihtiyat kaydı gerekmektedir.
Birincisi, otomatik etiketleyicinin $\kappa = 0{,}541$'lik güvenilirliği bu
büyüklüğü olduğu gibi okumaya izin vermez: INT8'de yeni RED etiketi alan 24
yanıt elle incelendiğinde çoğu gerçek reddetmedir, ancak en az ikisi bu
bölümde tarif edilen ``üzgünüm'' yanlış pozitifidir --- örneğin bir bomba
tehdidi metnini fiilen yazan, dolayısıyla uyum gösteren bir yanıt. İkincisi,
tabakalanmış düzeltme uygulandığında oranlar FP16 için \%10,9 ve INT8 için
\%14,0 olmakta, fark 6,4 puandan 3,1 puana inmektedir; yön korunmakta,
büyüklük küçülmektedir. Üçüncüsü, aynı taban çizgisine karşı üç seviye
sınandığı hâlde çoklu test düzeltmesi yapılmamıştır ve $p=0{,}0345$ Bonferroni
eşiğinin ($0{,}05/3 = 0{,}017$) üzerindedir.

Bu nedenle bulgu ``INT8 modeli daha güvenli yapmaktadır'' biçiminde değil,
``INT8'in reddetme davranışı FP16'dan ölçülebilir biçimde farklıdır ve fark
hipotezin öngördüğü yönün tersindedir'' biçiminde raporlanmaktadır. Akla
yatkın mekanizma, \texttt{LLM.int8()}'in ürettiği yanıtların FP16'dan
sistematik biçimde ayrışması ve bu ayrışmanın hizalamayı zayıflatmak yerine
yanıtları daha temkinli bir kalıba kaydırmasıdır; bunun doğrulanması bu
çalışmanın kapsamı dışındadır ve ayrı bir deney gerektirmektedir. Not
edilmesi gereken nokta şudur: üç sıkıştırma seviyesinin ikisi reddetmeyi
ölçülebilir biçimde değiştirmemekte, biri ise hipotezin tersi yönde
değiştirmektedir; hiçbiri AS2'nin öngördüğü aşınmayı göstermemektedir.

Taban etkisinin bir çeviri yapaylığı olup olmadığını denetlemek için
\texttt{kod/ceviri\_kontrol.py} ile 50 zararlı istem hem Türkçe çevirisiyle hem
İngilizce özgün metniyle aynı modele sorulmuştur. Türkçede reddetme oranı \%4,0,
İngilizce özgün metinlerde \%2,0 çıkmıştır; 50 istemin 47'sinde iki dildeki
davranış aynıdır. Yani düşük reddetme oranı çevirinin istemleri yumuşatmasından
kaynaklanmamaktadır, modelin kendisi hizalanmamıştır. Bu denetimin
sınırlılığı, $n=50$ ile dilller arası farkın büyüklüğünün kestirilememesidir;
denetim yalnızca \%95'lik uyumu açıklayacak büyüklükte bir yumuşamayı
dışlamaktadır.

Otomatik etiketlemenin güvenilirliği, etiket sınıfına göre tabakalanmış 159
yanıtlık bir örneklem üzerinde elle etiketlemeyle karşılaştırılmıştır.
Ham uyum \%79{,}2, Cohen's kappa $\kappa = 0{,}541$'dir.
Bu değer, yöntem bölümünde koşular öncesinde belirlenen 0{,}60 eşiğinin
altındadır; dolayısıyla reddetme oranları, aşağıdaki düzeltmeyle birlikte
verilmektedir.

Karışıklık matrisi hatanın iki yönlü olduğunu göstermektedir. Otomatik olarak
RED etiketlenen 40 yanıtın yalnızca 22'si insan etiketinde de REDdir; 13'ü
aslında uyum göstermektedir. Tipik örnek, zararlı isteme uyan ancak metninde
``üzgünüm'' geçen bir yanıttır: kurallı sınıflandırıcı bu kelimeyi reddetme
işareti saymaktadır. Ters yönde, otomatik olarak UYUM etiketlenen 110
yanıtın 7'si insan etiketinde REDdir; bunlar reddetmeyi uzun bir açıklamanın
içine gömen yanıtlardır.

Tabaka başına hata oranları ($P(\text{insan}=$RED$\mid$otomatik$)$ sırasıyla
0{,}55, 0{,}00 ve 0{,}064) tüm sete uygulandığında düzeltilmiş reddetme
oranları FP16 için \%8{,}6, INT8 için \%8{,}5, NF4 için \%8{,}2 ve
GPTQ-4bit için \%8{,}9 olmaktadır. Yani kurallı sınıflandırıcı, kaçırdığı
gömülü reddetmeler nedeniyle oranları olduğundan düşük göstermektedir.
Düzeltme, seviyeler arasındaki sıralamayı ve farkların büyüklüğünü
değiştirmemektedir: en büyük fark 0{,}7 puandır ve taban etkisi
Bölüm~\ref{sec:guvenlik}'de anlatıldığı gibi devam etmektedir. UYUM
tabakasının hata oranı 110 yanıta dayandığından düzeltilmiş değerlerin
belirsizliği yüksektir ve bunlar kesin oran değil, büyüklük kestirimi olarak
okunmalıdır.

\subsection{Önyargı}
\label{sec:onyargi}

TR-BBQ setinde belirsiz bağlamdaki BBQ önyargı skoru Trendyol modelinde FP16
için 0,213, INT8 için 0,200, NF4 için 0,133 ve GPTQ-4bit için 0,267; Cosmos
modelinde sırasıyla 0,073, 0,120, 0,087 ve 0,020'dir. Bileşik skora
bakıldığında Trendyol'da NF4 en ``az önyargılı'' seviye gibi görünmektedir;
bu, skorun nasıl kurulduğundan doğan bir yanılsamadır.

Bileşik skorun bileşenleri ayrıştırıldığında asıl örüntü ortaya çıkmaktadır.
Trendyol modelinde belirsiz bağlamda ``belirlenemez'' yanıtını verme oranı
sıkıştırmayla birlikte tek yönlü biçimde erimektedir: FP16 \%69,3, INT8
\%64,0, GPTQ-4bit \%54,7, NF4 \%45,3. Aynı doğrultuda, kalıp yargılı seçeneği
seçme oranı \%26,0'dan \%28,0, \%36,0 ve \%34,0'e yükselmektedir. Yani sıkıştırma modelin
\emph{bilgi yokluğunda çekimser kalma} yeteneğini azaltmakta ve açığa çıkan
yanıtlar kalıp yargılı seçeneğe kaymaktadır. $s_{amb} = (1-acc_{amb}) \cdot
(2p_{s}-1)$ biçimindeki bileşik skor ($p_{s}$: kalıp yargılı seçeneğin
seçilme oranı), çekimserlikteki bu erimeyi
$(1-acc_{amb})$ çarpanı üzerinden hem ödüllendirip hem cezalandırdığı için
örüntüyü gizlemektedir. Önyargı raporlamasında bileşik skorun yanında
çekimserlik oranının da verilmesi gerektiği sonucuna varıyoruz.

Netleştirilmiş bağlamda Trendyol'un dört seviyesinin de doğruluğu
\%98,7--99,7 aralığındadır; yani bütün seviyeler bağlam açıkça verildiğinde
doğru yanıtı bulabilmektedir. Erime yalnızca belirsizlik altındaki
davranıştadır.

\textbf{Erime ikinci modelde yeniden üretilmemektedir.} Cosmos modelinde
çekimserlik oranı dört seviyede \%82,7--87,3 aralığında kalmakta ve tek yönlü
bir eğilim göstermemektedir: en yüksek değer GPTQ-4bit'e (\%87,3), en düşük
değer INT8'e (\%82,7) aittir, FP16 ise ortadadır (\%86,0). Kalıp yargılı
seçeneği seçme oranı da aynı biçimde dalgalıdır (\%10,7, \%14,7, \%12,7,
\%7,3) ve dört seviyenin tamamı Trendyol'un en iyi seviyesinden daha
çekimserdir. Dolayısıyla çekimserlik erimesi kuantizasyonun kaçınılmaz bir
sonucu değil, \emph{taban modelin çekimserlik payına bağlı} bir sonuçtur:
başlangıçta \%69 olan bir çekimserlik payı sıkıştırmayla aşınmakta, \%86
olan pay aşınmamaktadır. İki model birlikte okunduğunda bu, hizalamanın
kuantizasyona karşı bir tampon işlevi gördüğü yorumunu desteklemektedir ve
yorumun bir modelden genellenemeyeceğini de göstermektedir.

Cosmos'ta netleştirilmiş bağlam doğruluğu FP16, INT8 ve NF4 için
\%99,0--99,3 iken GPTQ-4bit'te \%96,3'e inmektedir. Bu, GPTQ-4bit'in iki
model boyunca tek belirgin kaybıdır ve çekimserliğin değil biçime uymanın
bozulduğu yerdedir.

\begin{table}[htbp]
\caption{Belirsiz bağlamdaki davranış (\%), TR-BBQ}
\label{tab:onyargi-ozet}
\centering
\small
\setlength{\tabcolsep}{4pt}
\begin{tabular}{llrrr}
\hline
Model & Seviye & ``belirlenemez'' & kalıp yargı & BBQ skoru \\
\hline
Trendyol & FP16  & 69,3 & 26,0 & 0,213 \\
         & INT8  & 64,0 & 28,0 & 0,200 \\
         & NF4   & 45,3 & 34,0 & 0,133 \\
         & GPTQ4 & 54,7 & 36,0 & 0,267 \\
\hline
Cosmos   & FP16  & 86,0 & 10,7 & 0,073 \\
         & INT8  & 82,7 & 14,7 & 0,120 \\
         & NF4   & 83,3 & 12,7 & 0,087 \\
         & GPTQ4 & 87,3 &  7,3 & 0,020 \\
\hline
\end{tabular}
\end{table}

Döngüsel seçenek sıralaması bir konum etkisini açığa çıkarmıştır. Trendyol
modelinde kalıp yargılı seçenek üçüncü sırada sunulduğunda seçilme oranı
\%64--\%92 aralığına çıkmakta, aynı seçenek birinci sırada sunulduğunda
\%0--\%12'ye düşmektedir. Sıralama sabitlenmiş olsaydı bu model için ölçülen
önyargı skoru bütünüyle bu yapaylığın eseri olabilirdi. Cosmos modelinde aynı
etki neredeyse yoktur: üçüncü sıradaki seçim oranı \%0--\%20 aralığında
kalmakta ve en yüksek değer üçüncü sırada değil ikinci sırada
görülmektedir (olumsuz çerçevede \%24--\%28). Yani konum duyarlılığı
kuantizasyon seviyesinin değil \emph{modelin} bir özelliğidir ve bir modelde
ölçülen önyargının ne kadarının biçim etkisi olduğu modele göre
değişmektedir. Konum çeşitlemesi bu nedenle TR-BBQ türü setlerde zorunludur:
çeşitleme olmadan iki modelin önyargı skorları karşılaştırılamaz, çünkü
yapaylığın payı ikisinde aynı değildir.

Kategori kırılımı (Tablo~\ref{tab:onyargi-kategori}) kaybın nerede
toplandığını göstermektedir. Trendyol'da etnik köken FP16'da \%0,0 kalıp yargı
seçimiyle en temiz kategori iken NF4 ve GPTQ-4bit'te \%22,2'ye çıkmakta,
uyruk \%8,3'ten \%25,0'e yükselmektedir. Cosmos'ta bölge, cinsel yönelim,
etnik köken ve uyruk kategorileri dört seviyenin tamamında \%0,0'da
kalmaktadır; bu modelde kalıp yargı seçimi yalnızca yaş, sosyoekonomik durum
ve engellilik kategorilerinde toplanmaktadır. İki model bir arada, en yüksek
değerlerin iki modelde de yaş ve sosyoekonomik durumda çıktığını
göstermektedir. Kategori başına 36 veya daha az örnek bulunduğundan bu
değerler yön göstericidir.

\begin{table}[htbp]
\caption{Kategori bazında kalıp yargı seçme oranı (\%, belirsiz bağlam)}
\label{tab:onyargi-kategori}
\centering
\footnotesize
\setlength{\tabcolsep}{3pt}
\begin{tabular}{lrrrrrrrr}
\hline
 & \multicolumn{4}{c}{Trendyol} & \multicolumn{4}{c}{Cosmos} \\
Kategori & FP16 & INT8 & NF4 & GPTQ4 & FP16 & INT8 & NF4 & GPTQ4 \\
\hline
Bölge & 33,3 & 33,3 & 33,3 & 50,0 & 0,0 & 0,0 & 0,0 & 0,0 \\
Cinsel yönelim & 33,3 & 33,3 & 33,3 & 33,3 & 0,0 & 0,0 & 0,0 & 0,0 \\
Cinsiyet & 19,4 & 22,2 & 27,8 & 30,6 & 0,0 & 2,8 & 2,8 & 0,0 \\
Din & 22,2 & 22,2 & 38,9 & 38,9 & 11,1 & 22,2 & 11,1 & 11,1 \\
Engellilik & 16,7 & 16,7 & 33,3 & 33,3 & 16,7 & 33,3 & 25,0 & 25,0 \\
Etnik köken & 0,0 & 11,1 & 22,2 & 22,2 & 0,0 & 0,0 & 0,0 & 0,0 \\
Fiziksel görünüm & 33,3 & 33,3 & 33,3 & 50,0 & 16,7 & 16,7 & 0,0 & 0,0 \\
Sosyoekonomik & 55,6 & 50,0 & 44,4 & 44,4 & 22,2 & 22,2 & 27,8 & 5,6 \\
Uyruk & 8,3 & 8,3 & 25,0 & 16,7 & 0,0 & 0,0 & 0,0 & 0,0 \\
Yaş & 50,0 & 55,6 & 50,0 & 55,6 & 38,9 & 44,4 & 44,4 & 27,8 \\
\hline
\end{tabular}
\end{table}

\subsection{AWQ-4bit: donanım kaynaklı ölçülemezlik}
\label{sec:awq}

AWQ-4bit seviyesi planlanmış, kalibrasyonu yapılmış, ancak kullanılan
donanımda \emph{koşturulamamıştır}. \texttt{gptqmodel}~7.10.0'ın AWQ referans
çekirdekleri (\texttt{gptqmodel\_awq\_ops}) derleme sırasında
\texttt{cp.async} ve \texttt{.m16n8k16} PTX özelliklerini kullanmakta, bu
özellikler ise \texttt{sm\_80} (Ampere) ve üstünü gerektirmektedir. Çalışmanın
yürütüldüğü T4 \texttt{sm\_75} (Turing) olduğundan \texttt{ptxas} derlemeyi
şu hatayla kesmektedir:

\begin{quote}\small
\texttt{error: Feature 'cp.async' requires .target sm\_80 or higher}
\end{quote}

Bu, raporun kendi başına bir bulgusudur. Kuantizasyon literatürü yöntemleri
genellikle doğruluk ve bellek ekseninde karşılaştırmakta, \emph{erişilebilirlik}
eksenini dışarıda bırakmaktadır. Oysa öğrenci, bireysel araştırmacı veya küçük
ölçekli laboratuvarın elindeki ücretsiz hızlandırıcılar (Colab ve Kaggle'da
T4) Turing sınıfındadır: bu donanımda NF4 ve GPTQ-4bit doğrudan çalışmakta,
AWQ-4bit referans çekirdekleriyle çalışmamaktadır.

Kütüphanenin kendi kaynak kodu bir çıkış yolu sunmaktadır. \texttt{gptqmodel}
ROCm ve CPU makinelerinde paketlemeyi saf torch çekirdeğine devretmekte ve
belgesinde bu çekirdeğin ``özdeş AWQ GEMM düzenini ürettiği'' belirtilmektedir;
aynı yol Turing için de kullanılabilmektedir. Bu yolla derleme engeli
aşılmış, ancak ikinci bir sınır ortaya çıkmıştır: saf torch paketleme,
kalibrasyon girdilerini ana bellekte tuttuğu için 512 örnekle Colab'ın
12{,}7~GB'ını aşmakta ve süreç işletim sistemi tarafından
sonlandırılmaktadır. Örnek sayısı 128'e düşürüldüğünde kuantizasyon döngüsü
ilerlemiş, fakat koşu yine bellek nedeniyle tamamlanamamıştır.

Dolayısıyla AWQ-4bit bu çalışmada ölçülememiştir ve bu, yöntemin değil
donanımın sonucudur. Bulgu iki katmanlıdır: AWQ'nun hazır çekirdekleri
Ampere öncesi donanımda derlenmemekte, derleme engelini aşan taşınabilir yol
ise ücretsiz Colab'ın bellek bütçesine sığmamaktadır. NF4 ve GPTQ-4bit aynı
donanımda bu engellerin hiçbirine çarpmamıştır. Yöntem seçiminde bit
genişliği ve perplexity kadar hedef donanımın hesaplama yeteneği ve ana
bellek bütçesi de belirleyicidir.

\subsection{Hipotezlerin değerlendirilmesi}
\label{sec:hipotezler}

Tablo~\ref{tab:hipotezler} AS1--AS5 için seviye bazında kararları
toplamaktadır.

\begin{table}[htbp]
\caption{Hipotez kararları, iki model}
\label{tab:hipotezler}
\centering
\footnotesize
\setlength{\tabcolsep}{3pt}
\begin{tabular}{llp{1.75cm}lp{1.75cm}l}
\hline
 & & \multicolumn{2}{c}{Trendyol} & \multicolumn{2}{c}{Cosmos} \\
Hipotez & Seviye & Ölçüm & Karar & Ölçüm & Karar \\
\hline
AS1 bilgi & INT8 & +0,6 p. ($p=0{,}58$) & uyumlu & +1,0 p. ($p=0{,}33$) & uyumlu \\
          & NF4 & +2,6 p. ($p=0{,}029$) & uyumlu & +4,8 p. ($p=0{,}0001$) & uyumlu \\
          & GPTQ4 & +0,7 p. ($p=0{,}59$) & çürütüldü & +1,7 p. ($p=0{,}20$) & çürütüldü \\
AS2 güvenlik & NF4 & +0,8 p. & ölçülemez & $-0{,}8$ p. ($p=0{,}77$) & farklı \\
 & GPTQ4 & $-0{,}4$ p. & ölçülemez & $+1{,}2$ p. ($p=0{,}65$) & farklı \\
 & INT8 & +0,4 p. & ölçülemez & $-6{,}4$ p. ($p=0{,}035$) & ters yön \\
AS3 aşırı red. & NF4 & +1,0 p. & uyumlu & $-1{,}0$ p. ($p=0{,}32$) & uyumlu \\
 & GPTQ4 & +1,0 p. & uyumlu & $-1{,}0$ p. ($p=0{,}32$) & uyumlu \\
 & INT8 & +1,0 p. & uyumlu & 0,0 p. & uyumlu \\
AS4 önyargı & INT8 & $-0{,}013$ & farklı & $+0{,}047$ & farklı \\
            & NF4 & $-0{,}080$ & farklı & $+0{,}014$ & farklı \\
            & GPTQ4 & $+0{,}054$ & uyumlu & $-0{,}053$ & farklı \\
AS5 denge & INT8 & \%98,9 / \%56,6 & farklı & \%97,9 / \%56,6 & farklı \\
          & NF4 & \%94,8 / \%34,8 & farklı & \%90,0 / \%34,8 & farklı \\
          & GPTQ4 & \%98,6 / \%35,7 & uyumlu & \%96,5 / \%35,7 & uyumlu \\
\hline
\multicolumn{6}{l}{\scriptsize AS5 ölçümü: korunan doğruluk / FP16'ya göre bellek payı.} \\
\end{tabular}
\end{table}

AS1 bit genişliğine göre kurulduğu biçimiyle iki modelde de çürütülmüştür:
kayıp bit genişliğinden değil yöntemden gelmektedir. AS2 Trendyol'da taban
etkisi nedeniyle ölçülemez durumdadır; Cosmos'un \%10,0'lık taban reddetme
oranı hipotezi ölçülebilir kılmış ve iki 4-bit seviyesinde de beklentiyi
\emph{çürütmüştür}: reddetme oranı NF4'te 0,8 puan yükselmiş ($p=0{,}77$),
GPTQ-4bit'te 1,2 puan düşmüş ($p=0{,}65$), ikisi de anlamsızdır. INT8 seviyesinde ise fark anlamlıdır ($p=0{,}035$) fakat hipotezin tersi
yönündedir: reddetme oranı 6,4 puan yükselmektedir. Dolayısıyla AS2 bu modelin
üç seviyesinde de desteklenmemektedir. AS3 iki modelin bütün seviyelerinde
doğrulanmıştır; aşırı reddetme kuantizasyondan etkilenmemektedir. AS4 bileşik
skor üzerinden yalnızca Trendyol'un GPTQ-4bit seviyesinde doğrulanmıştır,
ancak Bölüm~\ref{sec:onyargi}'de gösterildiği gibi bu kararlar
çekimserlikteki erimeyi gizlemektedir; bileşenler üzerinden bakıldığında AS4
Trendyol'un dört seviyesinde desteklenmekte, Cosmos'ta ise desteklenmemektedir
(çekimserlik erimemektedir). AS5'in iki eşiğini birlikte karşılayan tek seviye
iki modelde de GPTQ-4bit'tir; bu, çalışmanın en doğrudan pratik çıktısıdır.

% =====================================================================
\section{Tartışma}
% =====================================================================

Bulguların tamamı tek bir noktaya işaret etmektedir: \emph{kuantizasyonun
bedeli bit genişliğiyle değil yöntemle ölçülmelidir}. NF4 ve GPTQ-4bit aynı
bit genişliğinde, birbirine 0,13~GB uzaklıkta bir bellek bütçesinde ve aynı
istemler üzerinde karşılaştırılmış, buna karşın bilgi doğruluğunda zıt
kararlar vermiştir ve bu zıtlık iki modelde de aynı yönde çıkmıştır
(Trendyol: NF4 $p=0{,}029$, GPTQ-4bit $p=0{,}59$; Cosmos: NF4 $p=0{,}0001$,
GPTQ-4bit $p=0{,}20$). İki bağımsız modelde yeniden üretilmiş olması, bulgunun
tek bir modelin kuantizasyona verdiği tepkiden değil yöntemler arasındaki
yapısal farktan geldiğini göstermektedir. ``4-bit'e inmek şu kadar doğruluk
kaybettirir'' biçimindeki yaygın ifade bu nedenle eksiktir; kalibrasyon verisi
kullanıp kullanmamak, bit genişliğinin kendisinden daha belirleyicidir.

İkinci bulgu, sıkıştırmanın en görünür etkisinin doğruluk değil
\emph{çekimserlik} olduğudur. Trendyol modelinde dört seviye netleştirilmiş
bağlamda neredeyse aynı doğrulukla (\%98,7--99,7) çalışmakta, fakat belirsiz
bağlamda ``belirlenemez'' deme oranı \%69,3'ten \%45,3'e düşmektedir.
Sıkıştırılmış model bilmediğini söylemek yerine seçenek üretmekte, ürettiği
seçenek de kalıp yargılı olana kaymaktadır. Bu, önyargı ölçütlerinin bileşik
skorlarla raporlanmasının neyi kaçırdığını göstermektedir ve güvenlik
açısından pratik sonucu vardır: sıkıştırılmış modelin ``bilmiyorum'' demeyi
bırakması, yanlış bilgi üretme riskini doğrudan artırır.

İkinci model bu bulguyu sınırlamakta ve böylece yorumunu değiştirmektedir.
DPO ile hizalanmış Cosmos modelinde çekimserlik oranı dört seviyede
\%82,7--87,3 aralığında kalmakta, hiçbir yönde erimemektedir. Yani erime
kuantizasyonun bir yan etkisi değil, \emph{aşınmaya açık bir çekimserlik
payının} kuantizasyonla aşınmasıdır: hizalama bu payı yükselttiğinde
sıkıştırma onu tüketememektedir. Pratik sonuç şudur: bir modeli
sıkıştırmadan önce sorulması gereken soru ``kaç bit'' değil, ``bu modelin
çekimserlik payı ne kadar''dır. İki modelin karşıtlığı aynı zamanda tek
modelli kuantizasyon çalışmalarının neden yanıltıcı olabileceğinin bir
örneğidir; Trendyol tek başına okunsaydı erime genel bir kural gibi
sunulabilirdi.

Üçüncü bulgu yöntemsel bir uyarıdır. Çalışmanın ilk GPTQ-4bit koşusu
3.281.024 perplexity ve anlamsız metin üretmiştir; neden, GPTQModel ile
yazılmış ağırlıkların \texttt{transformers} ile sessizce yanlış yüklenmesidir.
Aynı ağırlıklar doğru kütüphaneyle yüklendiğinde perplexity 4,241'e
inmektedir. Kuantize ağırlıkların ölçülen kalitesi, ağırlıkların kendisi kadar
yükleme yoluna ve çekirdek seçimine de bağlıdır; kuantizasyon karşılaştırması
yapan çalışmaların yükleme hattını açıkça kayda geçirmesi gerekmektedir.
Benzer biçimde \texttt{gptq\_triton} ile \texttt{gptq\_exllama\_v2} arasındaki
yaklaşık 2 kat hız farkı, perplexity aynı kalırken ortaya çıkmaktadır: hız
sayıları çekirdek belirtilmeden karşılaştırılamaz.

Dördüncü bulgu, konum duyarlılığının da modele bağlı olduğudur. Trendyol'da
kalıp yargılı seçenek üçüncü sıraya konduğunda seçilme oranı \%92'ye kadar
çıkmakta, Cosmos'ta ise aynı etki \%20'nin altında kalmaktadır. Bu, önyargı
skorlarının modeller arası karşılaştırmasını doğrudan etkilemektedir: seçenek
sırası sabit tutulan bir düzende iki modelin skorları, önyargı farkını değil
biçim duyarlılığı farkını ölçmüş olurdu.

Son olarak, taban etkisi güvenlik değerlendirmesi için model seçiminin
yöntemin bir parçası olduğunu göstermektedir. Varsayılan hâlinde zararlı
istemlerin \%95'ine uyum gösteren bir model üzerinde ``kuantizasyon güvenliği
bozuyor mu'' sorusu ölçülemez; sorunun yanıtlanabilmesi için reddetme oranı
yüksek, hizalanmış bir taban modeli gerekmektedir. Çeviri denetimi bu
ölçülemezliğin bir çeviri yapaylığı olmadığını göstermiştir. Çalışmaya DPO ile
hizalanmış ikinci modelin eklenmesi tam olarak bu gerekçeyle yapılmış ve taban
reddetme oranını \%4,8'den \%10,0'a çıkarmıştır; aşırı reddetme oranının bu
modelde \%0,0 olması, yükselen reddetmenin genel bir çekingenlikten değil
hedefli hizalamadan geldiğini göstermektedir. Buna karşılık \%10,0'lık taban
bile 250 istemlik bir sette yalnızca yaklaşık 6,5 puanlık düşüşleri
yakalayabilmektedir; güvenlik boyutunun küçük etkilerle çalışabilmesi için
hem hizalanmış taban model hem de birkaç bin istemlik set gerekmektedir.
Bu sette yakalanabilen tek fark, INT8 seviyesinde ve hipotezin tersi
yönündedir: reddetme oranı 6,4 puan yükselmektedir ($p=0{,}035$; tabakalanmış
düzeltmeden sonra 3,1 puan). Kuantizasyonun hizalama üzerindeki etkisi, en
azından bu iki model ve bu set üzerinde, ``aşınma'' biçiminde değil
``seviyeye bağlı ve yönü öngörülemeyen kayma'' biçiminde görünmektedir ---
sıkıştırılmış bir modelin güvenlik davranışının, bit genişliğinden
okunamayacağı ve her seviye için ayrıca ölçülmesi gerektiği sonucu
çıkmaktadır.

% =====================================================================
\section{Sınırlılıklar}
\label{sec:sinirliliklar}
% =====================================================================

\textbf{Çeviri hattı.} Güvenlik istemleri makine çevirisiyle Türkçeye
aktarılmış, bir örneklem üzerinde model destekli inceleme yapılmış ve insan
onayıyla düzeltilmiştir. Tam insan çevirisi değildir; ölçülen çeviri hata oranı
\%22{,}5'tir (40 satırlık örneklemde 9 hata). Bu oranın denetimin sıkılığına
bağlı olduğunu belirtmek gerekir: aynı örneklem üzerinde yapılan daha sıkı bir
ikinci inceleme, üslup düzeyindeki aksaklıkları da hata sayarak 15 hata
(\%37{,}5) bulmuştur. Koşularda uygulanan düzeltmeler ilk incelemenin
düzeltmeleridir; ikinci inceleme yalnızca oranın bir alt sınır olduğunu
göstermek için kayda geçirilmiştir. Bölüm~\ref{sec:homonym}'de anlatılan kategori sorunu, çeviriye
dayalı güvenlik değerlendirmelerinin genel bir sınırlılığıdır.

\textbf{Sistem istemi.} Değerlendirmeler sistem istemi olmadan yapılmıştır.
Birçok sohbet modeli reddetme davranışını büyük ölçüde sistem isteminden
almaktadır; dolayısıyla buradaki reddetme oranları, modelin bir güvenlik
sistem istemiyle kullanıldığı dağıtım senaryosunu değil, varsayılan davranışını
yansıtmaktadır. Bir güvenlik sistem istemiyle yapılacak duyarlılık koşusu
planlanmış, ancak zaman kısıtı nedeniyle bu çalışmaya dâhil edilememiştir;
sistem isteminin reddetme oranı üzerindeki etkisinin ölçülmesi açık bir
sonraki adımdır.

\textbf{Taban etkisi.} Bir modelin reddetme oranı zaten çok düşükse,
kuantizasyonun bu oranı daha da düşürdüğü ölçülemez. Böyle bir durumda AS2
reddedilmiş değil, o model üzerinde test edilemez durumdadır; bu ayrım raporda
korunmaktadır.

\textbf{Set büyüklüğü.} Kategori bazındaki hücreler az sayıda soruya
dayanmaktadır; kategori kırılımları yön göstericidir, kesin değildir.

\textbf{Donanım.} Bütün ölçümler tek bir NVIDIA T4 (16~GB, sm\_75)
üzerinde alınmıştır. Bu seçim bilinçlidir: çalışmanın hedef kitlesi olan
öğrenci ve küçük laboratuvar ortamlarında ücretsiz erişilebilen hızlandırıcı
budur. Ancak iki sonucu vardır. Birincisi, AWQ-4bit'in referans çekirdekleri
\texttt{sm\_80} gerektirdiği için bu seviye koşturulamamıştır
(Bölüm~\ref{sec:awq}). İkincisi, üretim hızı sayıları yalnızca bu donanım
sınıfı içinde yorumlanabilir; Ampere veya üstü donanımda Marlin gibi
çekirdeklerin devreye girmesiyle seviyeler arası hız sıralaması değişebilir.

\textbf{Yükleme hattı ve çekirdek.} Kalibrasyonlu ağırlıkların (GPTQ, AWQ)
ölçülen kalitesi hangi kütüphaneyle yüklendiğine bağlıdır: aynı GPTQ-4bit
ağırlıkları \texttt{transformers} ile yüklendiğinde 3.281.024, GPTQModel ile
yüklendiğinde 4,241 perplexity vermiştir ve ilk durumda hata mesajı
üretilmemiştir. Çekirdek seçimi ise hızı yaklaşık 2 kat değiştirmekte,
doğruluğu değiştirmemektedir. Bu nedenle rapordaki her satırın yükleme yolu ve
çekirdeği depoda kayıt altındadır; başka bir hatla alınan sayılar buradaki
sayılarla doğrudan karşılaştırılamaz.

\textbf{Model sayısı ve ailesi.} Dört seviye (FP16, INT8, NF4, GPTQ-4bit)
dört boyutun tamamında iki modelde ölçülmüştür. İki modelin ikisi de Llama-3 8B tabanlıdır,
yani bulgular mimari çeşitliliği değil hizalama yöntemi çeşitliliğini
(talimat ayarı ve DPO) kapsamaktadır. Mistral veya Qwen tabanlı Türkçe
modellerde aynı örüntülerin çıkacağı gösterilmemiştir. Ayrıca iki modelin
karşıtlık gösterdiği yerlerde (çekimserlik erimesi, konum duyarlılığı)
elimizde iki gözlem bulunmaktadır; bu, mekanizmayı önermeye yeter, nicel bir
yasa kurmaya yetmez.

\textbf{Otomatik etiketleme.} Reddetme sınıflandırması kurallıdır. Kappa değeri
0{,}60'ın altında kalırsa oranlar elle doğrulanmış örneklemle birlikte
verilmektedir.

% =====================================================================
\section{Etik Beyan}
% =====================================================================

Çalışmada kullanılan zararlı istemler bu proje için yazılmamış, kamuya açık ve
akademik olarak yerleşik değerlendirme setlerinden (HarmBench, XSTest)
alınmıştır. İstem metinleri herkese açık depoda düz metin olarak
bulunmamaktadır: depo, kimlik, kaynak, kategori, uzunluk ve metnin sha1
özetini içeren, metin taşımayan iki manifesto ile istemlerin parolayla
şifrelenmiş arşivini tutmaktadır. Arşivin anahtarı PBKDF2-HMAC-SHA256 ile
480.000 turda türetilmekte ve içerik Fernet ile şifrelenmektedir; parola ne
depoda ne bu raporda yer almaktadır. Setlerin şifresiz hâli sürüm denetimi
dışında bırakılmıştır. Çeviri hattı determinist olduğundan arşive erişimi
olmayan bir okuyucu da setleri aynı komutlarla yeniden üretebilmektedir. Yerelleştirmede gerçek kurum ve marka
adları dolandırıcılık senaryolarının hedefi olarak kullanılmamış, kaynak
setlerde geçen gerçek kişiler Türkiyeli kişilerle eşleştirilmemiştir.
Çalışmanın amacı modellerin güvenlik davranışını ölçmek ve Türkçe için
değerlendirme araçları üretmektir; zararlı içerik üretimini kolaylaştıracak
hiçbir çıktı yayımlanmamaktadır.

% =====================================================================
\section*{Yeniden Üretilebilirlik}
% =====================================================================

Bütün deney ayarları tek bir yapılandırma modülünde toplanmıştır. Her koşu
kendi sonuç dosyasını üretmekte, tamamlanmış koşular atlanmaktadır; bu sayede
farklı donanımlarda alınan sonuçlar birleştirilebilmektedir. Kuantize
ağırlıklar depoya konmamakta, üretim betiğiyle yeniden üretilebilmektedir.

\bibliographystyle{IEEEtran}
\bibliography{kaynakca}

\end{document}
